\section{Related decision targets and the cost of disjoint pairing}
\label{app:pairing_efficiency}

\paragraph{Stage reweighting and raw component requirements differ.}
The priority-standardized net benefit of \citet{mccoy2026psnb} combines
stage-conditional effects using prespecified weights. Consider the same
success-first, success-eligible cost rule as in Section~\ref{sec:estimands}.
Let $B$ always succeed at cost 2; let $A$ fail with probability
$\epsilon=0.1$ and otherwise succeed at cost 1. Success is the first
tier. The cost tier is reached exactly when $A$ succeeds, with
probability $1-\epsilon=0.9$, and every reached comparison favors $A$.
Thus the stage-conditional effects are $(-0.1,1)$, ordinary net benefit
is $-0.1+0.9=0.8$, and priority-standardized net benefit under weights
$(0.8,0.2)$ is $-0.08+0.2=0.12$. Both composite effects are positive,
while the raw success contrast $-0.1$ fails a margin of $0.03$.
All stage reach probabilities are positive, and costs are decisive under
the 5\% tolerance. This algebra illustrates different population
requirements. It does not show a miscalibrated test, compare operating
characteristics, or claim novelty for the lower-tier-dominance concern.

\paragraph{A standard projection identity at the same execution budget.}
An efficient all-pairs estimator can reduce variance by reusing complete
arm records. The following identity quantifies the fixed-prefix price of
discarding those comparisons. It is the usual two-sample projection
decomposition, not a new U-statistic result. It applies to the independent
stationary arm law used by our primary score simulation; adaptive,
nonexchangeable, or task-coupled records require a separate argument.

Let $A_1,\ldots,A_N\overset{\mathrm{iid}}\sim F_A$ and
$B_1,\ldots,B_N\overset{\mathrm{iid}}\sim F_B$, with independent arms.
Both procedures use $2N$ executed episodes: $N$ records per arm.
For a square-integrable comparison $h$, put
\[
 \theta=\E h(A,B),\quad
 a(A)=\E\{h(A,B)\mid A\}-\theta,\quad
 b(B)=\E\{h(A,B)\mid B\}-\theta,
\]
and $r(A,B)=h(A,B)-\theta-a(A)-b(B)$. Define
$\sigma_A^2=\E a^2$, $\sigma_B^2=\E b^2$, and
$\sigma_R^2=\E r^2$. Consider
\[
 D_N=\frac1N\sum_{i=1}^N h(A_i,B_i),\qquad
 U_N=\frac1{N^2}\sum_{i,j=1}^N h(A_i,B_j).
\]
\begin{proposition}[Exact fixed-prefix variance comparison]
Under the independent-arm law above,
\[
 \operatorname{Var}(D_N)
   =\frac{\sigma_A^2+\sigma_B^2+\sigma_R^2}{N},\qquad
 \operatorname{Var}(U_N)
   =\frac{\sigma_A^2+\sigma_B^2}{N}+\frac{\sigma_R^2}{N^2}.
\]
Consequently, $\operatorname{Var}(D_N)-\operatorname{Var}(U_N)
=(N-1)\sigma_R^2/N^2\geq0$.
\end{proposition}
\begin{proof}
Independence gives $\E a=\E b=0$ and
$\E(r\mid A)=\E(r\mid B)=0$. Thus $a$, $b$, and $r$ are orthogonal
in $L^2$, so $\operatorname{Var}(h)=\sigma_A^2+\sigma_B^2+\sigma_R^2$.
The $N$ disjoint scores are independent, proving the first formula.
For the all-pairs statistic,
\[
 U_N-\theta=\frac1N\sum_i a(A_i)+\frac1N\sum_j b(B_j)
                    +\frac1{N^2}\sum_{i,j}r(A_i,B_j).
\]
The three displayed terms are pairwise uncorrelated. Two distinct
residual terms are independent when they share no index. When they
share only an $A$ index, condition on that $A$ record and use independent
$B$ records and $\E(r\mid A)=0$; their covariance is again zero.
The same argument applies to a shared $B$ index. Only the $N^2$
identical-index residual variances remain, giving
$N^2\sigma_R^2/N^4=\sigma_R^2/N^2$. Adding the projection variances
proves the second formula and the stated difference.
\end{proof}

If $\sigma_A^2+\sigma_B^2>0$, the variance ratio is
\[
 \frac{\operatorname{Var}(D_N)}{\operatorname{Var}(U_N)}
 =\frac{\sigma_A^2+\sigma_B^2+\sigma_R^2}
        {\sigma_A^2+\sigma_B^2+\sigma_R^2/N}.
\]
Its large-$N$ limit depends on the score law. If the first-order
projections vanish but $\sigma_R^2>0$, the ratio is $N$; if all three
variances vanish, both estimators are constant and their ratio is
undefined. At $N=1$ the estimators coincide.

For a component kernel $h_g(A,B)=g(A)-g(B)$, the residual is exactly
zero, and $D_N=U_N=\bar g_A-\bar g_B$ for every realized sample.
This applies to the raw success and compliance gates in our simulation.
Reusing all cross-arm comparisons can improve the hierarchical score's
point-estimation variance but provides no additional component-mean
observations on these same records. Different confidence boundaries may
still have different behavior.

These fixed-prefix identities do not compare anytime error control,
conjunction stopping time, delayed observation policies, or computation.
All-pairs scores are dependent and cannot be inserted as $N^2$
independent observations into the disjoint-pair betting process.
The sequential U-statistic literature
\citep{zhang2024sequential,bergemann2026group,cai2026ustatistics}
addresses that dependence under its respective assumptions. No
empirical superiority claim over those procedures follows from this
identity or from our current reference experiments.
